\label{section:Conclusion}
The primary objective of this paper was to compare the compilation results of C and Rust using two scripts with the Lazart tool. Our findings reveal that Rust compiles the code differently than C, causing LLVM to interpret the code in a distinct manner. Consequently, certain attacks may evade detection by the Lazart tool, resulting in the loss of some vulnerabilities. This discrepancy is illustrated by the tables generated from the C code [\ref{tab1:memcmps_C}, \ref{tab3:verifypin_C}] and those from the Rust code [\ref{tab2:memcmps_Rust}, \ref{tab4:verifypin_Rust}].

Upon analyzing the control-flow graph, we noticed that Rust may optimize certain blocks of code with special instructions such as \texttt{select} or \texttt{switch}.
While these optimizations enhance performance, they can obfuscate the control flow and reduce the effectiveness of static analysis tools like Lazart.

The second objective of this paper was to explore methods for preventing such optimizations during compilation. Our findings indicate that the inline(never) flag is crucial to ensure Lazart behaves as intended. This flag forces Rust to generate a call instruction for the function, preserving it in its entirety and making it fully analyzable. Additionally, our findings suggest that renaming blocks in the LLVM IR could prevent further optimizations of these blocks, thereby improving the accuracy of static analysis.

Another promising area for further investigation is the use of global variables as function parameters to prevent the loss of additional vulnerabilities. We hypothesize that the compiler cannot optimize global variables passed as parameters due to a lack of assumptions about their usage. This approach could provide a practical solution to retain critical information for static analysis tools.

These findings highlight the challenges and opportunities in using static analysis tools like Lazart across different programming languages and compilation environments. While our results demonstrate the impact of compiler optimizations on vulnerability detection, they also open new avenues for improving static analysis tools by adapting them to modern compilation practices.

In future work, we aim to extend Wolverine to handle the special instructions generated by Rust and create more benchmarks for Lazart in Rust.

%\end{multicols}